\documentclass[12pt]{article}
\usepackage[utf8]{inputenc}
\usepackage[T1]{fontenc}
\usepackage{hyperref}
\title{Aranama-Tamique Language Resources for K-12 Curriculum Development}
\author{Compiled by Claude Sonnet 5.5 (claude-sonnet-5-5) for Andr\'as Kornai}
\date{}
\begin{document}
\maketitle

\section{Introduction}
Aranama-Tamique (ISO 639-3: xrt; Glottocode aran1265), also called Araname, Haname or Tamique, is an extinct, unclassified language of southern Texas. The only known record is a single word and a two-word phrase that Albert Gatschet took down in 1884 from a Tonkawa man called ``Old Simon'', who said they came from a language he called ``Hanama''. The words given are \emph{himiy\'ana} ``water'' and \emph{himia'na tsa'yi} ``Give me water!''. Goddard (1979) is reported to have shown that this is not enough evidence to link the language to any documented language of southern Texas. A heritage and revitalization curriculum in the language is therefore not possible with the available record, and no revitalization programme was found.

\section{Primary Educational Resources}
None found.

\section{Linguistic and Reference Materials}
\begin{itemize}
\item Wikipedia, Aranama language: \url{https://en.wikipedia.org/wiki/Aranama_language}
\item Wikipedia, Coahuiltecan languages (a possible, unproven affiliation): \url{https://en.wikipedia.org/wiki/Coahuiltecan_languages}
\item Swanton 1940 source record in the ASJP database: \url{https://asjp.clld.org/sources/1953}
\end{itemize}

\section{Religious and Community Resources}
No community organisation of Aranama descendants was found. Historical background on the Aranama people, who lived at the missions at Refugio and San Antonio:
\begin{itemize}
\item Wikipedia, Aranama people: \url{https://en.wikipedia.org/wiki/Aranama_people}
\item Texas Beyond History, Native peoples of the Spanish mission period: \url{https://texasbeyondhistory.net/espiritu/native.html}
\item National Park Service, San Antonio Missions (Campbell study): \url{https://www.nps.gov/parkhistory/online_books/saan/campbell/sec3.htm}
\end{itemize}

\section{Digital and Audio Resources}
None found. There are no recordings.

\section{Curriculum Development Context}
Only a few words survive, so no language instruction can be built. A K-12 heritage unit could cover the history of the Aranama people and the Texas missions using the historical sources above, and could teach how little can be known about a language from such a record. Any such unit should involve Texas Native American and mission-descendant communities, none of which was identified in these searches.

\end{document}